%==============================================================================
%  MetaLab Suite --- Theoretical Framework
%  Trajectory-Based Metaheuristics for Wind Farm Layout Optimization (WFLO)
%  and the Dynamic Drone Routing Problem (DDRP)
%
%  Compile with:  pdflatex theoretical_framework.tex   (run twice for the ToC)
%==============================================================================
\documentclass[11pt,a4paper]{article}

\usepackage[T1]{fontenc}
\usepackage[utf8]{inputenc}
\usepackage{lmodern}
\usepackage[margin=2.5cm]{geometry}
\usepackage{amsmath,amssymb,amsthm,mathtools}
\usepackage{booktabs}
\usepackage{array}
\usepackage{algorithm}
\usepackage{algpseudocode}
\usepackage{xcolor}
\usepackage{caption}
\usepackage{microtype}
\usepackage[hidelinks]{hyperref}

%------------------------------------------------------------------------------
%  Mathematical shortcuts and theorem-like environments
%------------------------------------------------------------------------------
\newtheorem{definition}{Definition}
\newtheorem{proposition}{Proposition}
\newtheorem{remark}{Remark}

\DeclareMathOperator*{\argmax}{arg\,max}
\DeclareMathOperator*{\argmin}{arg\,min}
\newcommand{\R}{\mathbb{R}}
\newcommand{\E}{\mathbb{E}}
\newcommand{\Prob}{\mathbb{P}}
\newcommand{\norm}[1]{\left\lVert #1 \right\rVert}
\newcommand{\abs}[1]{\left\lvert #1 \right\rvert}

\title{%
  \textbf{Theoretical Framework for Trajectory-Based Metaheuristics}\\[3mm]
  \large Wind Farm Layout Optimization and the Dynamic Drone Routing Problem}
\author{%
  MetaLab Suite\\[1mm]
  \small Department of Operations Research}
\date{\today}

%==============================================================================
\begin{document}
\maketitle

\begin{abstract}
\noindent
This document establishes the theoretical foundations of the MetaLab Suite, an
educational repository devoted to \emph{trajectory-based metaheuristics} for
combinatorial and continuous optimisation. We develop, in full mathematical
detail, two flagship problems. \textbf{Part~I} addresses the
\emph{Wind Farm Layout Optimization} (WFLO) problem: given a terrain and a wind
resource, where should turbines be placed so as to maximise the expected power
output while mitigating the aerodynamic \emph{wake effect}? We derive the
objective function $P_{\text{total}}$, state the topographical and engineering
constraints, and present a simplified \emph{Jensen} wake model, before arguing
why the resulting landscape is strongly multimodal and therefore calls for
trajectory-based search. \textbf{Part~II} addresses the \emph{Dynamic Drone
Routing Problem} (DDRP) in the spirit of Dorling et al.~(2016): a fleet of
unmanned aerial vehicles serves delivery orders that are revealed in real time,
subject to a strict payload capacity $Q$ and a battery budget whose consumption
scales with both payload and travelled distance. We formulate the objective
(minimise distance plus unserved-order penalties) and the associated
constraints, extend the model to the dynamic setting, and close with the design
of trajectory metaheuristics for both problems.
\end{abstract}

\tableofcontents
\bigskip

%==============================================================================
\section{Introduction}
%==============================================================================

\subsection{Scope and motivation}

Many real-world planning problems are \emph{black-box}, \emph{constrained} and
\emph{multimodal}: the objective cannot be differentiated analytically, feasible
regions are riddled with holes, and the landscape contains a combinatorial
profusion of local optima. Two representative examples motivate this document:
the siting of wind turbines on a complex terrain (Part~I) and the routing of
delivery drones under real-time demand (Part~II).

Trajectory-based metaheuristics---also called \emph{single-solution} or
\emph{local-search} metaheuristics---are particularly well suited to such
problems. Rather than maintaining a population, they evolve a single incumbent
solution $x^{0}, x^{1}, \dots, x^{T}$ by iteratively exploring a
\emph{neighbourhood} $\mathcal{N}(x)$ and applying an acceptance rule that may
temporarily worsen the objective so as to escape local optima.

\subsection{Trajectory-based metaheuristics: common skeleton}

Formally, given a finite or continuous \emph{search space} $\mathcal{X}$, an
objective $f:\mathcal{X}\to\R$ (to be minimised or maximised), and a
neighbourhood operator $\mathcal{N}:\mathcal{X}\to 2^{\mathcal{X}}$, a generic
trajectory method proceeds as in Algorithm~\ref{alg:generic}.

\begin{algorithm}[H]
\caption{Generic trajectory-based metaheuristic}
\label{alg:generic}
\begin{algorithmic}[1]
\State Initialise $x \leftarrow x^{0} \in \mathcal{X}$ (e.g.\ at random or greedily)
\State $x^{\star} \leftarrow x$ \Comment{best solution found so far}
\While{stopping criterion not met}
    \State Generate a candidate $x' \in \mathcal{N}(x)$
    \If{$\textsc{Accept}(x, x', \text{state})$}
        \State $x \leftarrow x'$
    \EndIf
    \If{$f(x')$ improves $f(x^{\star})$}
        \State $x^{\star} \leftarrow x'$
    \EndIf
    \State Update the search \emph{state} (temperature, tabu tenure, \dots)
\EndWhile
\State \Return $x^{\star}$
\end{algorithmic}
\end{algorithm}

Three instantiations are used throughout the suite and differ only in the
acceptance rule:
\begin{itemize}
  \item \textbf{Simulated annealing (SA)} accepts a worsening move with
        probability $\exp(-\Delta / T)$, where $\Delta>0$ is the objective
        degradation and $T$ a slowly decreasing temperature
        \cite{kirkpatrick1983}.
  \item \textbf{Tabu search (TS)} always moves to the best admissible neighbour
        but forbids recently visited configurations through a short-term memory
        (the \emph{tabu list}), with an aspiration criterion to override the
        tabu status on sufficiently good moves \cite{glover1989}.
  \item \textbf{Iterated local search (ILS)} alternates a local-search descent
        with a stochastic \emph{perturbation} of the incumbent, accepting the
        result by a simple criterion \cite{lourenco2003}.
\end{itemize}

A crucial property for Part~II is that trajectory methods are
\emph{anytime}: they may be interrupted at any moment and still return the best
solution found so far, which is exactly what a real-time dispatcher needs.

\subsection{Notation}

\begin{table}[H]
\centering
\caption{Global notation used in both parts.}
\label{tab:notation}
\begin{tabular}{@{}ll@{}}
\toprule
Symbol & Meaning \\
\midrule
$\mathcal{X}$            & search space (problem-dependent encoding) \\
$f(\cdot)$               & objective function of the metaheuristic \\
$\mathcal{N}(x)$         & neighbourhood of solution $x$ \\
$x^{\star}$              & best solution found during a run \\
$T$                      & temperature (SA) \\
$\mathcal{T}$            & tabu list (TS) \\
$N$                      & number of turbines (Part~I) \\
$Q$                      & payload capacity of a drone (Part~II) \\
$B$                      & battery energy budget of a drone (Part~II) \\
\bottomrule
\end{tabular}
\end{table}

%==============================================================================
\section{Part I --- Wind Farm Layout Optimization (WFLO)}
%==============================================================================

\subsection{Problem description}

A wind farm occupies a domain $\Omega \subset \R^{2}$ (the available land,
possibly irregular and with an elevation profile). We must choose the positions
\[
  X = \bigl(x_{1}, x_{2}, \dots, x_{N}\bigr), \qquad
  x_{j} = \bigl(x_{j}^{(1)}, x_{j}^{(2)}\bigr) \in \Omega,
\]
of $N$ identical turbines so that the farm's electrical production is as large
as possible over a reference year. The fundamental difficulty is the
\emph{wake effect}: an upwind turbine extracts kinetic energy from the flow and
leaves behind it a region of reduced wind speed and increased turbulence. A
turbine standing in that wake therefore produces less than it would in
free-stream conditions. Pushing turbines apart reduces wake losses but consumes
land and increases collector-system and civil costs; clustering them together
saves land but intensifies interference. The optimal design must balance these
two effects, which is what makes the problem non-trivial.

\subsection{The wake effect}

Let the wind travel in direction $\theta$ (the bearing it blows towards), with
unit vector $\mathrm{n}=(\cos\theta,\sin\theta)$ and its normal
$\mathrm{n}^{\perp}=(-\sin\theta,\cos\theta)$. For two turbines $i$ and $j$,
define the \emph{along-wind} and \emph{cross-wind} separations
\begin{equation}
  x_{ij} = (x_{j}-x_{i})\cdot\mathrm{n},
  \qquad
  y_{ij} = \bigl\lvert (x_{j}-x_{i})\cdot\mathrm{n}^{\perp} \bigr\rvert .
  \label{eq:separations}
\end{equation}
A deficit is created behind turbine $i$ only if $x_{ij}>0$ (turbine $j$ is
downwind of $i$) and if $j$ lies inside the expanding wake.

\subsection{Simplified Jensen wake model}

The Jensen (or Kati\'c--Jensen) model \cite{jensen1983,katic1986} assumes a
\emph{top-hat} velocity profile: inside the wake the wind speed is uniform, and
the wake expands linearly with downstream distance at a rate governed by the
decay constant $k$. Its radius at downstream distance $x$ from a rotor of
diameter $D$ is
\begin{equation}
  R_{w}(x) = \frac{D}{2} + k\,x .
  \label{eq:jensen-radius}
\end{equation}
When turbine $j$ is partially or fully immersed in the wake of $i$, the
fractional velocity deficit is
\begin{equation}
  \delta_{ij}
  = \left( 1-\sqrt{1-C_{T}} \right)
    \left( \frac{D}{D + 2k\,x_{ij}} \right)^{2},
  \qquad
  \delta_{ij}=0 \ \text{if } x_{ij}\le 0 \ \text{or } y_{ij} > R_{w}(x_{ij}),
  \label{eq:jensen-deficit}
\end{equation}
where $C_{T}\in(0,1)$ is the thrust coefficient of the turbine. The factor
$1-\sqrt{1-C_{T}}$ equals twice the axial induction factor $a$, so
Eq.~\eqref{eq:jensen-deficit} has a clear physical interpretation: the maximum
deficit immediately behind the rotor is $2a$, and it decays as the wake spreads.

For a given wind scenario the total deficit at turbine $j$ is obtained by
superposing the contributions of all upwind turbines, most commonly in
\emph{root-sum-square} form,
\begin{equation}
  U_{j} = U_{\infty}\left( 1-\sqrt{\sum_{i \in \mathcal{I}_{j}} \delta_{ij}^{2}} \right)_{+},
  \qquad
  \mathcal{I}_{j} = \bigl\{\, i : x_{ij}>0 \ \wedge\ y_{ij} \le R_{w}(x_{ij}) \,\bigr\},
  \label{eq:effective-speed}
\end{equation}
with $U_{\infty}$ the free-stream speed and $(\cdot)_{+}=\max(\cdot,0)$ to keep
the effective speed non-negative. Linear superposition
$U_{j}=U_{\infty}(1-\sum_{i}\delta_{ij})$ is a common alternative; both are
supported by the reference implementation. The Jensen model is deliberately
simple. It neglects turbulence anisotropy, ground effects and yaw, yet it
reproduces the qualitative behaviour that dominates the optimisation and keeps
the objective cheap enough for the tens of thousands of evaluations a
metaheuristic requires.

\subsection{Turbine power curve}

The electrical power produced by a single turbine as a function of wind speed
$u$ is modelled by the standard piecewise \emph{power curve}. With cut-in speed
$u_{\text{in}}$, rated speed $u_{r}$, cut-out speed $u_{\text{out}}$ and rated
power $P_{r}$,
\begin{equation}
  P(u) =
  \begin{cases}
    0, & u < u_{\text{in}} \ \text{or} \ u \ge u_{\text{out}},\\[1mm]
    P_{r}\left(\dfrac{u-u_{\text{in}}}{u_{r}-u_{\text{in}}}\right)^{3},
      & u_{\text{in}} \le u < u_{r},\\[2mm]
    P_{r}, & u_{r} \le u < u_{\text{out}}.
  \end{cases}
  \label{eq:power-curve}
\end{equation}
The cubic regime reflects the fact that the kinetic power of the wind scales
with the cube of its speed, $P_{\text{wind}}=\tfrac12\rho A u^{3}$, where
$\rho$ is air density and $A$ the rotor swept area.

\subsection{Wind resource: the discretised wind rose}

The wind climate is summarised by a \emph{wind rose}: a finite set of $S$
scenarios $\{(\theta_{s}, U_{s}^{\infty}, p_{s})\}_{s=1}^{S}$, where
$\theta_{s}$ is a direction, $U_{s}^{\infty}$ a representative speed and
$p_{s}\in[0,1]$ its relative frequency, with $\sum_{s=1}^{S}p_{s}=1$. The rose is
typically obtained by binning historical anemometry data in direction and speed.

\subsection{Objective function: expected power \texorpdfstring{$P_{\text{total}}$}{P\_total}}

For scenario $s$, applying \eqref{eq:effective-speed} and \eqref{eq:power-curve}
to every turbine gives the farm's instantaneous power
\begin{equation}
  P_{\text{total}}^{(s)}(X)
  = \sum_{j=1}^{N} P\bigl(U_{j}^{(s)}(X)\bigr),
  \label{eq:power-scenario}
\end{equation}
and the \emph{expected} power output of the layout is the probability-weighted
average over the rose,
\begin{equation}
  \boxed{\;
  \bar{P}(X)
  = \sum_{s=1}^{S} p_{s}\, P_{\text{total}}^{(s)}(X)
  = \sum_{s=1}^{S} p_{s} \sum_{j=1}^{N}
      P\!\left( U_{\infty,s}
      \Bigl(1-\sqrt{\textstyle\sum_{i\in\mathcal{I}_{j}} \delta_{ij}^{2}}\Bigr)_{+} \right).
  \;}
  \label{eq:objective-wflo}
\end{equation}
Multiplying by the number of hours in a year, $T_{\text{year}}=8760\ \text{h}$,
yields the \emph{Annual Energy Production} $\text{AEP}(X)=T_{\text{year}}\,
\bar{P}(X)$. Because turbines may be assigned a hub height $H_{j}=h(x_{j})+H_{0}$
that follows the terrain elevation $h$, the effective speeds in
\eqref{eq:objective-wflo} can be evaluated at the local hub height, although we
keep the notation compact here. The WFLO problem is then
\begin{equation}
  \text{(WFLO)} \qquad
  \max_{X \in \mathcal{F}} \ \bar{P}(X),
  \label{eq:wflo}
\end{equation}
where $\mathcal{F}$ is the feasible set defined next.

\subsection{Topographical and engineering constraints}

The feasible set $\mathcal{F}=\mathcal{F}_{1}\cap\mathcal{F}_{2}\cap
\mathcal{F}_{3}\cap\mathcal{F}_{4}\cap\mathcal{F}_{5}$ encodes the physical and
regulatory limits.

\begin{description}
  \item[($\mathcal{F}_{1}$) Land availability / boundary clearance.] Every
  turbine, including its rotor, must lie inside the buildable domain $\Omega$:
  \begin{equation}
    \operatorname{dist}\!\bigl(x_{j}, \partial\Omega\bigr) \ \ge \ R = D/2,
    \qquad j=1,\dots,N .
    \label{eq:c-boundary}
  \end{equation}
  On real terrain $\Omega$ excludes water bodies, settlements, roads and other
  protected areas.

  \item[($\mathcal{F}_{2}$) Exclusion zones.] A collection of forbidden
  polygons $\{Z_{m}\}_{m=1}^{M}$ (buildings, radar cones, archaeological sites)
  must be avoided:
  \begin{equation}
    x_{j} \notin \bigcup_{m=1}^{M} Z_{m},
    \qquad j=1,\dots,N .
    \label{eq:c-exclusion}
  \end{equation}

  \item[($\mathcal{F}_{3}$) Minimum inter-turbine spacing.] To limit wake
  interference and mechanical loading, turbines must keep a centre-to-centre
  distance of at least a multiple $\gamma$ of the rotor diameter:
  \begin{equation}
    \norm{x_{i}-x_{j}} \ \ge \ s_{\min} = \gamma D,
    \qquad 1 \le i < j \le N,
    \label{eq:c-spacing}
  \end{equation}
  with typical $\gamma \in [3,8]$.
  Eq.~\eqref{eq:c-spacing} makes $\mathcal{F}$ \emph{non-convex}: it removes an
  open neighbourhood around each turbine, effectively carving holes in the
  domain.

  \item[($\mathcal{F}_{4}$) Terrain and slope.] Turbine foundations require
  limited ground slope and, for road access, acceptable gradients between
  positions. A common surrogate bounds the elevation difference relative to the
  horizontal separation,
  \begin{equation}
    \frac{\abs{h(x_{i})-h(x_{j})}}{\norm{x_{i}-x_{j}}}
      \ \le \ \sigma_{\max},
    \qquad 1\le i<j\le N ,
    \label{eq:c-slope}
  \end{equation}
  and hub heights are constrained to a permissible band
  $H_{\min}\le h(x_{j})+H_{0}\le H_{\max}$.

  \item[($\mathcal{F}_{5}$) Discrete/engineering constraints.] Optional
  restrictions such as a maximum number of turbines along a line, alignment
  rules, or a grid-snapping of positions to an accessible lattice.
\end{description}

\subsection{Why the landscape is multimodal, and why trajectory search fits}

The WFLO objective \eqref{eq:objective-wflo} is, for several independent
reasons, a pathological landscape for gradient-based methods.

\begin{proposition}[Structural difficulties]
\label{prop:wflo-hard}
The WFLO objective $\bar{P}$ over $\mathcal{F}$ exhibits the following
properties.
\begin{enumerate}
  \item \textbf{Non-convexity and disconnectedness.} The spacing constraints
        \eqref{eq:c-spacing} render the feasible set non-convex; it is a union
        of many disjoint regions in $\R^{2N}$.
  \item \textbf{Non-separability.} Because of the wake terms $\delta_{ij}$,
        the contribution of turbine $j$ depends on the positions of all
        turbines upwind of it. The objective does not decompose into a sum of
        per-turbine terms, so turbines cannot be optimised independently.
  \item \textbf{Nonsmoothness.} The indicator condition
        $y_{ij}\le R_{w}(x_{ij})$ in \eqref{eq:jensen-deficit} introduces a
        discontinuity in the first derivative of $\bar{P}$ whenever a turbine
        enters or leaves a wake cone; likewise the $(\cdot)_{+}$ and the
        piecewise power curve are non-differentiable at their kinks.
  \item \textbf{Combinatorial multimodality.} Each of the
        $\binom{N}{2}$ wake interactions contributes a local basin, so the
        number of local optima grows rapidly with $N$. Good layouts typically
        arrange turbines in staggered patterns aligned with the dominant wind
        direction, and distinct wind-rose structures give rise to competing,
        mutually exclusive configurations.
\end{enumerate}
\end{proposition}

Classical nonlinear programming solvers stall on the first local optimum and
cannot cross the spacing barriers of (F3). This is precisely the regime in
which \emph{trajectory-based} metaheuristics excel: they treat $\bar{P}$ as a
black box, they navigate non-convex feasible sets through clever encodings and
penalties, and their acceptance rules (Metropolis, tabu tenure, perturbation)
let them escape the numerous local optima catalogued in
Proposition~\ref{prop:wflo-hard}.

\paragraph{Encoding and neighbourhood.}
Because the natural variables are continuous, trajectory methods operate
directly on the coordinate vector $X\in\R^{2N}$. A standard neighbourhood
perturbs a single turbine,
\begin{equation}
  \mathcal{N}_{\text{move}}(X)
  = \bigl\{\, X + e_{j}\otimes \xi \ : \ 1\le j\le N,\ \xi\sim\mathcal{N}(0,\sigma^{2}I_{2}) \,\bigr\},
  \label{eq:wflo-neighbourhood}
\end{equation}
optionally combined with \emph{swap} moves (exchange two turbines) and
\emph{block} moves (translate a subset of turbines). Infeasible candidates are
handled either by rejection, by a repair operator, or by a penalty
$\lambda\sum_{i<j}[s_{\min}-\norm{x_{i}-x_{j}}]_{+}$.

\paragraph{Reference algorithm.}
Algorithm~\ref{alg:sa-wflo} instantiates the skeleton with a Metropolis
acceptance and a geometric cooling schedule $T_{t+1}=\alpha T_{t}$; this is the
reference metaheuristic against which the baseline hill climber of the
repository is compared.

\begin{algorithm}[H]
\caption{Simulated annealing for WFLO (maximisation of $\bar{P}$)}
\label{alg:sa-wflo}
\begin{algorithmic}[1]
\State $X \leftarrow$ \Call{RandomFeasibleLayout}{$N,\Omega$}
\State $X^{\star} \leftarrow X$;\quad $T \leftarrow T_{0}$
\While{$T > T_{\text{end}}$ \textbf{and} budget remains}
  \Repeat
    \State Draw $X' \in \mathcal{N}_{\text{move}}(X)$ and repair/clip to $\mathcal{F}$
    \State $\Delta \leftarrow \bar{P}(X') - \bar{P}(X)$
    \If{$\Delta > 0$ \textbf{or} $\operatorname{rand}() < \exp(\Delta/T)$}
      \State $X \leftarrow X'$
      \If{$\bar{P}(X) > \bar{P}(X^{\star})$}
        \State $X^{\star} \leftarrow X$
      \EndIf
    \EndIf
  \Until{$L$ iterations completed}
  \State $T \leftarrow \alpha T$
\EndWhile
\State \Return $X^{\star}$
\end{algorithmic}
\end{algorithm}

%==============================================================================
\section{Part II --- Dynamic Drone Routing Problem (DDRP)}
%==============================================================================

\subsection{Motivation and related work}

Unmanned aerial vehicles (drones) are increasingly used for last-mile parcel
delivery. Dorling et al.~\cite{dorling2016} introduced and studied the
\emph{Vehicle Routing Problem for drone delivery}, in which a fleet of
quad-rotor drones, launched from a depot, serves a set of customers under hard
limits on payload weight and battery energy. Their model demonstrates that the
cost of a delivery is governed jointly by \emph{distance} and by
\emph{payload-dependent energy} consumption, so the routing and the loading
decisions are inseparable.

In the \emph{Dynamic} Drone Routing Problem (DDRP) studied here, delivery
orders are not known in advance: they are revealed in \emph{real time} over a
planning horizon. The dispatcher must continuously reassign drones, implying a
trade-off between exploiting the current plan and re-optimising as new
information arrives. Dynamic routing of this kind generalises the static
vehicle-routing problem and is NP-hard \cite{toth2014}; its online nature makes
it even more challenging \cite{psaraftis1995}.

\subsection{Sets, parameters and decision variables}

\begin{table}[H]
\centering
\caption{DDRP data and decision variables.}
\label{tab:ddrp}
\begin{tabular}{@{}ll@{}}
\toprule
Symbol & Meaning \\
\midrule
$\mathcal{K}=\{1,\dots,K\}$      & set of drones \\
$\mathcal{I}=\{1,\dots,n\}$      & set of delivery orders (customers) \\
$0$                              & the depot (\emph{departure/return node}) \\
$V=\{0\}\cup\mathcal{I}$         & node set \\
$d_{ij}$                         & travel distance between nodes $i,j$ \\
$w_{i}$                          & payload (parcel weight) of order $i$ \\
$Q$                              & payload capacity of a drone \\
$B$                              & usable battery energy of a drone \\
$a$                              & empty-drone energy consumption per unit distance \\
$b$                              & payload energy sensitivity per unit distance \\
$v$                              & cruise speed (distance per unit time) \\
$\tau$                           & service (hover/landing) time per order \\
$\lambda$                        & penalty for an unserved order \\
\addlinespace
$x_{ijk}\in\{0,1\}$              & $1$ if drone $k$ travels directly from $i$ to $j$ \\
$y_{ik}\in\{0,1\}$               & $1$ if drone $k$ serves order $i$ \\
$q_{ijk}\ge 0$                   & payload carried by drone $k$ on arc $(i,j)$ \\
$t_{ik}\ge 0$                    & (auxiliary) position of node $i$ in route $k$ \\
$z_{k}\in\{0,1\}$                & $1$ if drone $k$ is used \\
\bottomrule
\end{tabular}
\end{table}

\subsection{Payload-dependent battery model}

A multi-rotor's power requirement increases with its total mass, hence with the
payload it carries. Following the linearisation used by Dorling et al.\
\cite{dorling2016}, the \emph{energy consumed per unit distance} by a drone
carrying a payload $q$ is
\begin{equation}
  e(q) = a + b\,q, \qquad a \ge 0,\ b \ge 0 ,
  \label{eq:energy-rate}
\end{equation}
so that the energy spent traversing arc $(i,j)$ with payload $q$ is
$e(q)\,d_{ij}$. The two constants capture the empty-drone consumption $a$ and
the extra consumption caused by carrying mass ($b\,q$). This dependency is what
makes DDRP different from a pure distance-minimisation problem: a heavy parcel
carried far consumes disproportionately more battery.

\subsection{Objective function}

Let $y_{ik}=\sum_{j\in V} x_{ijk}$ indicate whether drone $k$ serves order $i$.
The problem trades off three quantities: total travelled distance, total energy
consumed, and the number of orders left unserved. The primary objective
requested here---\emph{minimise distance and unserved penalties}---reads
\begin{equation}
  \boxed{\;
  \min \ Z(X)
  = \underbrace{\sum_{k\in\mathcal{K}} \sum_{i\in V} \sum_{j\in V}
      d_{ij}\, x_{ijk}}_{\text{total travelled distance}}
  \;+\; \lambda
    \underbrace{\sum_{i\in\mathcal{I}} \Bigl(1-\sum_{k\in\mathcal{K}} y_{ik}\Bigr)}_{\text{unserved orders}}
  \;}
  \label{eq:objective-ddrp}
\end{equation}
where $\lambda\gg \max_{ij}d_{ij}$ is chosen large enough that serving an order
is always preferable to abandoning it, yet finite so that genuinely
unreachable orders under the battery budget are penalised rather than rendering
the instance infeasible. If energy is to be priced explicitly (as in
\cite{dorling2016}), the objective is augmented with
$c_{e}\sum_{k}\sum_{i,j} (a + b\,q_{ijk})\,d_{ij}\,x_{ijk}$, where $c_{e}$ is
the cost per unit of energy; we keep the two-term form
\eqref{eq:objective-ddrp} as the canonical teaching objective.

\subsection{Constraints}

The following constraints are \emph{strict}: no plan may violate them.

\paragraph{(C1) Unique assignment and unserved accounting.}
Each order is served by at most one drone,
\begin{equation}
  \sum_{k\in\mathcal{K}} y_{ik} \ \le \ 1,
  \qquad i \in \mathcal{I}.
  \label{eq:ddrp-assign}
\end{equation}

\paragraph{(C2) Flow conservation and route structure.}
For every served node the in-degree equals the out-degree, and each used drone
leaves and re-enters the depot,
\begin{equation}
  \sum_{j\in V} x_{ijk} = \sum_{j\in V} x_{jik} = y_{ik},
  \quad i\in\mathcal{I},\ k\in\mathcal{K};
  \qquad
  \sum_{j\in\mathcal{I}} x_{0jk} = \sum_{i\in\mathcal{I}} x_{i0k} = z_{k},
  \quad k\in\mathcal{K}.
  \label{eq:ddrp-flow}
\end{equation}

\paragraph{(C3) Payload capacity.}
The payload flowing into a customer equals the payload flowing out plus the
parcel delivered there; the depot loads each drone with exactly the parcels it
will deliver:
\begin{align}
  \sum_{j\in V} q_{ijk} - \sum_{j\in V} q_{jik} &= -\,w_{i}\,y_{ik},
    & i &\in \mathcal{I},\ k\in\mathcal{K},
    \label{eq:ddrp-load-balance}\\
  \sum_{j\in V} q_{0jk} &= \sum_{i\in\mathcal{I}} w_{i}\,y_{ik},
    & k &\in\mathcal{K},
    \label{eq:ddrp-depot-load}
\end{align}
and the capacity limit holds on every arc and for the whole load,
\begin{equation}
  0 \ \le \ q_{ijk} \ \le \ Q\,x_{ijk},
  \quad i,j\in V,\ k\in\mathcal{K};
  \qquad
  \sum_{i\in\mathcal{I}} w_{i}\,y_{ik} \ \le \ Q\,z_{k},
  \quad k\in\mathcal{K}.
  \label{eq:ddrp-capacity}
\end{equation}

\paragraph{(C4) Battery budget.}
The energy consumed along the whole route, with consumption scaling with
payload \eqref{eq:energy-rate} and distance, must not exceed the usable battery
energy,
\begin{equation}
  \sum_{i\in V}\sum_{j\in V} \bigl(a + b\,q_{ijk}\bigr)\, d_{ij}\, x_{ijk}
  \ \le \ B\, z_{k},
  \qquad k\in\mathcal{K}.
  \label{eq:ddrp-battery}
\end{equation}
In practice one reserves a safety margin $B_{\text{res}}$ for wind, hovering and
landing, replacing $B$ by $B-B_{\text{res}}$.

\paragraph{(C5) Subtour elimination.}
Auxiliary order variables with a big-$M$ bound forbid disconnected subtours,
\begin{equation}
  t_{ik} - t_{jk} + n\,x_{ijk} \ \le \ n-1,
  \qquad i,j\in\mathcal{I},\ i\neq j,\ k\in\mathcal{K}.
  \label{eq:ddrp-mtz}
\end{equation}

\paragraph{(C6) Release times (dynamic information).}
An order $i$ with release time $r_{i}$ may be served only once it is known. Let
$t^{a}_{ik}$ denote the arrival time at order $i$. Then
\begin{equation}
  y_{ik} = 1 \ \Rightarrow \ t^{a}_{ik} \ \ge \ r_{i},
  \label{eq:ddrp-release}
\end{equation}
which we linearise via $t^{a}_{ik} \ge r_{i} - M(1-y_{ik})$. Optional time
windows $[\,\underline{e}_{i},\overline{e}_{i}\,]$ extend this to
$\underline{e}_{i} \le t^{a}_{ik} \le \overline{e}_{i}$ when $y_{ik}=1$.

\subsection{The dynamic setting}

In DDRP the order set $\mathcal{I}$ is not known a priori. Instead orders
arrive according to a stochastic process---typically a non-homogeneous Poisson
process with rate $\nu(t)$ over a horizon $[0,T]$---and each order becomes
\emph{actionable} at its release time $r_{i}$. The dispatcher observes the
state
\[
  S_{t} = \bigl(\text{positions and remaining batteries of all drones},\
          \text{served orders},\ \text{pending orders}\bigr)
\]
and selects an action: dispatch a drone, insert a customer into an existing
route, or idle. A \emph{policy} $\pi$ maps states to actions; the Dynamic DDRP
seeks
\begin{equation}
  \min_{\pi} \ \E\Bigl[ \ \sum_{t=0}^{T}
      \Bigl( \textstyle\sum_{k}\sum_{i,j} d_{ij}\,x_{ijk}^{t}
             \;+\; \lambda\, U_{t} \Bigr) \ \Bigr],
  \label{eq:ddrp-dynamic}
\end{equation}
where $U_{t}$ is the number of still-unserved, already-released orders at epoch
$t$. The expectation is taken over the (unknown) arrival process.

A tractable and widely used approach is \emph{rolling-horizon
re-optimisation}: at each decision epoch $\Delta$ the dispatcher solves a
\emph{static} instance consisting of all pending orders plus new arrivals, with
the drones' current positions, remaining payload and remaining battery as
initial conditions, and commits to the first portion of the plan before
re-optimising at the next epoch. This decomposition turns the online problem
into a sequence of offline DDRP instances, each solvable by the trajectory
metaheuristics of Section~\ref{sec:ddrp-meta}.

\subsection{Trajectory metaheuristics for the DDRP}
\label{sec:ddrp-meta}

\paragraph{Encoding.}
A convenient single-solution encoding is the \emph{giant tour}: a permutation
of the orders, split optimally (by dynamic programming) into $K$ feasible
drone routes. The split respects capacity~\eqref{eq:ddrp-capacity} and
battery~\eqref{eq:ddrp-battery} and is recomputed after every move. Unserved
orders are simply those the split cannot place, and are charged $\lambda$ in
\eqref{eq:objective-ddrp}.

\paragraph{Neighbourhoods.}
Classic VRP moves operate on the permutation:
\begin{itemize}
  \item \textbf{relocate} --- remove one order and reinsert it elsewhere;
  \item \textbf{exchange} --- swap two orders;
  \item \textbf{Or-opt} --- move a chain of $2$ or $3$ consecutive orders;
  \item \textbf{2-opt} --- reverse a segment, useful within a route.
\end{itemize}

\begin{algorithm}[H]
\caption{Tabu search for the (static) DDRP}
\label{alg:ts-ddrp}
\begin{algorithmic}[1]
\State $s \leftarrow$ \Call{CheapestInsertion}{pending orders}
\State $s^{\star} \leftarrow s$;\quad $\mathcal{T} \leftarrow \emptyset$
\While{stopping criterion not met}
  \State $s' \leftarrow \argmin_{\hat{s}\in\mathcal{N}(s)} \{ Z(\hat{s}) : \text{move}(\hat{s}) \notin \mathcal{T} \}$
  \If{$s'$ is tabu but $Z(s') < Z(s^{\star})$} \Comment{aspiration}
      \State release $s'$ from $\mathcal{T}$
  \EndIf
  \State $s \leftarrow s'$;\quad insert the reverse move into $\mathcal{T}$ with tenure $\theta$
  \If{$Z(s) < Z(s^{\star})$}
      \State $s^{\star} \leftarrow s$
  \EndIf
\EndWhile
\State \Return $s^{\star}$
\end{algorithmic}
\end{algorithm}

\paragraph{Why trajectory methods suit the dynamic setting.}
Trajectory methods are the natural choice for DDRP because (i) they are
\emph{anytime}, so the dispatcher can halt them when a new order arrives and
use the best solution available; (ii) they \emph{warm-start} trivially, by
seeding the search with the previous plan so that real-time re-optimisation
inherits past effort; and (iii) the giant-tour representation admits cheap,
incremental evaluation of relocate/exchange moves. This closes the conceptual
circle opened in the introduction: the same skeleton of
Algorithm~\ref{alg:generic} serves both the continuous, multiplication-of-wakes
landscape of Part~I and the discrete, payload-and-battery-constrained routing
of Part~II.

%==============================================================================
\section{Synthesis}
%==============================================================================

Table~\ref{tab:synthesis} contrasts the two problems. Despite their different
mathematical natures---one continuous, one combinatorial---both are
constrained, black-box and multimodal, and both are attacked in the MetaLab
Suite with the same family of trajectory-based metaheuristics, instantiated
with problem-specific encodings and neighbourhoods.

\begin{table}[H]
\centering
\caption{Comparison of the two problems.}
\label{tab:synthesis}
\begin{tabular}{@{}p{3.1cm}p{5.3cm}p{5.3cm}@{}}
\toprule
Aspect & WFLO (Part~I) & DDRP (Part~II) \\
\midrule
Decision            & continuous turbine coordinates $X\in\R^{2N}$
                    & discrete routes and assignments \\
Objective           & maximise $\bar{P}(X)$ / AEP
                    & minimise distance $+\ \lambda\,(\text{unserved})$ \\
Key physics/economics & Jensen wake deficit; wind rose
                    & payload-dependent battery $e(q)=a+bq$ \\
Hard constraints    & boundary, exclusion, spacing, slope
                    & capacity $Q$, battery $B$, precedence, release times \\
Source of multimodality & wake coupling; non-convex feasible set
                    & NP-hard routing; online arrivals \\
Trajectory method   & simulated annealing / tabu on coordinates
                    & tabu search / ILS on the giant tour \\
Special property    & derivative-free, nonsmooth objective
                    & \emph{anytime}, warm-startable for re-optimisation \\
\bottomrule
\end{tabular}
\end{table}

%==============================================================================
\begin{thebibliography}{99}

\bibitem{dorling2016}
K.~Dorling, J.~Heinrichs, G.~G.~Messier, and S.~Magierowski,
``Vehicle routing problems for drone delivery,''
\emph{IEEE Transactions on Systems, Man, and Cybernetics: Systems},
vol.~47, no.~1, pp.~70--85, 2017 (available online 2016).

\bibitem{jensen1983}
N.~O.~Jensen,
``A note on wind generator interaction,''
Ris\o{} National Laboratory, Roskilde, Denmark, Tech.\ Rep.\ Ris\o-M-2411, 1983.

\bibitem{katic1986}
I.~Kati\'c, J.~H\o{}jstrup, and N.~O.~Jensen,
``A simple model for cluster efficiency,''
in \emph{Proc.\ European Wind Energy Association Conference and Exhibition},
Rome, Italy, 1986, pp.~407--410.

\bibitem{kirkpatrick1983}
S.~Kirkpatrick, C.~D.~Gelatt, and M.~P.~Vecchi,
``Optimization by simulated annealing,''
\emph{Science}, vol.~220, no.~4598, pp.~671--680, 1983.

\bibitem{glover1989}
F.~Glover,
``Tabu search---Part I,''
\emph{ORSA Journal on Computing}, vol.~1, no.~3, pp.~190--206, 1989.

\bibitem{lourenco2003}
H.~R.~Louren\c{c}o, O.~C.~Martin, and T.~St\"utzle,
``Iterated local search,''
in \emph{Handbook of Metaheuristics}, F.~Glover and G.~A.~Kochenberger, Eds.
Boston, MA: Springer, 2003, pp.~320--353.

\bibitem{mosetti1994}
G.~Mosetti, C.~Poloni, and B.~Diviacco,
``Optimization of wind turbine positioning in large wind farms by means of a
genetic algorithm,''
\emph{Journal of Wind Engineering and Industrial Aerodynamics},
vol.~51, no.~1, pp.~105--116, 1994.

\bibitem{grady2005}
S.~A.~Grady, M.~Y.~Hussaini, and M.~M.~Abdullah,
``Placement of wind turbines using genetic algorithms,''
\emph{Renewable Energy}, vol.~30, no.~2, pp.~259--270, 2005.

\bibitem{toth2014}
P.~Toth and D.~Vigo, Eds.,
\emph{Vehicle Routing: Problems, Methods, and Applications}, 2nd ed.
Philadelphia, PA: SIAM, 2014.

\bibitem{psaraftis1995}
H.~N.~Psaraftis,
``Dynamic vehicle routing: Status and prospects,''
\emph{Annals of Operations Research}, vol.~61, no.~1, pp.~143--164, 1995.

\bibitem{bent2004}
R.~W.~Bent and P.~Van Hentenryck,
``Scenario-based planning for partially dynamic vehicle routing with stochastic
customers,''
\emph{Operations Research}, vol.~52, no.~6, pp.~977--987, 2004.

\bibitem{ropke2006}
S.~Ropke and D.~Pisinger,
``An adaptive large neighborhood search heuristic for the pickup and delivery
problem with time windows,''
\emph{Transportation Science}, vol.~40, no.~4, pp.~455--472, 2006.

\end{thebibliography}

\end{document}
